\chapter{Oficio: cómo armar un modelo macro que funcione}
\label{cap:tips}

Este capítulo y el siguiente no tratan del algoritmo sino del oficio. Son las
cosas que uno aprende perdiendo semanas y que nadie escribe en los artículos.
Están ordenadas por el momento en que aparecen al trabajar.

\section{Antes de escribir una ecuación}

\paragraph{Escribe primero la pregunta, después el modelo} Suena obvio y casi
nadie lo hace. La pregunta determina qué heterogeneidad hace falta. Si la
pregunta es ``cuánto cae el producto'', probablemente baste con un agente
representativo. Si es ``quién absorbe la caída'', hace falta heterogeneidad, y
hace falta \emph{en la dimensión de la pregunta}. Un modelo con heterogeneidad
en una dimensión irrelevante es más caro y no es más informativo.

\paragraph{Construye por capas, y resuelve cada capa} La estrategia de empezar
con un agente representativo, luego agregar una fuente de heterogeneidad, luego
otra, es correcta --- y conviene \emph{resolver numéricamente} cada capa antes de
agregar la siguiente, no solo escribirla. Cada capa nueva debe reproducir la
anterior cuando se apaga el parámetro que la activa. Ese es el mejor test de
regresión que existe.

\paragraph{Decide qué es exógeno antes de calibrar} La tentación de endogeneizar
todo es fuerte y casi siempre es un error en una tesis de pregrado. Cada variable
endógena adicional es una incógnita, un objetivo, una condición de identificación
y una fuente de problemas de determinación.

\section{Al escribir el modelo}

\paragraph{Convenciones de tiempo: elige una y escríbela} La fuente número uno de
errores en modelos dinámicos es la ambigüedad sobre qué significa un subíndice.
¿$k_t$ es el capital elegido en $t$ o el disponible en $t$? Las dos convenciones
son defendibles; mezclarlas no. Conviene escribir en un comentario del código,
en la primera línea, cuál se usa.

\paragraph{Los rezagos van en los bloques simples} Como se vio en el capítulo
\ref{cap:extensiones}, un bloque heterogéneo no admite rezagos del insumo. La
disciplina es: el bloque heterogéneo recibe solo variables contemporáneas; los
rezagos y adelantos se construyen fuera.

\paragraph{Normaliza algo} En casi todo modelo hay una indeterminación de escala.
Conviene fijarla explícitamente (la masa de agentes igual a uno, el producto de
estado estacionario igual a uno, la productividad media igual a uno) y
verificarlo numéricamente.

\paragraph{Usa la ley de Walras como verificación, no como ecuación} Si se
imponen todas las condiciones de vaciado, una es redundante y $\mathbf{H}_U$ es
singular. La práctica correcta: omitir una condición del sistema y \emph{usarla
como prueba}. Si el residuo de la condición omitida no es numéricamente cero, hay
un error de contabilidad.

\section{Al elegir incógnitas y objetivos}

Esta es la decisión que más afecta la velocidad y la estabilidad, y la que menos
se discute.

\begin{receta}[Elegir el conjunto de incógnitas]
\begin{enumerate}[leftmargin=1.4em,itemsep=2pt]
\item Lista todas las variables endógenas.
\item Marca las que puedas despejar explícitamente en función de otras. Esas
salen del sistema: van a bloques simples.
\item De las que quedan, identifica las que generan \emph{ciclos} en el grafo.
Esas tienen que ser incógnitas, con su condición de consistencia como objetivo.
\item Verifica que el número de incógnitas iguale el de objetivos, y que cada
objetivo sea sensible a al menos una incógnita.
\end{enumerate}
\end{receta}

Un DAG eficiente puede ser diez veces más rápido que uno ``plano'' donde todas
las variables son incógnitas. En el modelo HANK de dos activos de \abrs, el DAG
eficiente tiene 3 incógnitas frente a 18 del plano, y resuelve siete veces más
rápido.

\begin{trampa}
Un objetivo insensible a toda incógnita produce una fila de ceros en
$\mathbf{H}_U$ y una matriz singular. Suele pasar al copiar y pegar: uno declara
un objetivo que en realidad ya está satisfecho por construcción. El síntoma es un
error de álgebra lineal; la causa está en el planteo, no en el código numérico.
\end{trampa}

\section{Al calibrar}

\paragraph{Calibra el estado estacionario con objetivos, no con parámetros} En
lugar de fijar $\beta=0.98$ y ver qué sale, fija la relación
capital--producto que quieres y resuelve por $\beta$. Es más transparente, más
fácil de defender y hace el modelo comparable con otros.

\paragraph{Verifica que el estado estacionario sea interior} Este punto es
específico de este método y crítico. Antes de calcular ningún jacobiano, hay que
verificar que ninguna función $\min$, $\max$ o restricción agregada esté
\emph{exactamente} en su quiebre, y conviene reportar a qué distancia está. Si
$\min(1, P/Q)$ se evalúa en $P/Q = 0.98$, el modelo es formalmente diferenciable
pero un choque modesto lo cruza, y la linealización deja de ser informativa.

\paragraph{Cuidado con las condiciones de crecimiento} En modelos con acumulación
y retornos aproximadamente lineales, la existencia de una distribución
estacionaria no trivial depende de una condición del tipo $\beta R < 1$ (o, con
salida de agentes, $\beta\,\varsigma\,R<1$). Si la condición se viola, la
distribución diverge y el código no converge o, peor, converge a una distribución
con masa apilada en el extremo de la grilla. Hay que \emph{mirar} la masa en el
último punto de la grilla: si no es del orden de $10^{-10}$ o menos, la grilla
está restringiendo el modelo.

\section{Al depurar}

\begin{receta}[Orden de depuración cuando algo no cuadra]
\begin{enumerate}[leftmargin=1.4em,itemsep=2pt]
\item ¿Converge el estado estacionario a $10^{-14}$? ¿Suma uno la distribución?
¿Se cumple el presupuesto agregado \emph{exactamente}?
\item ¿Coinciden los jacobianos del bloque con el método directo a $T=30$?
\item ¿Se cumplen las identidades analíticas que conoces (respuestas en impacto
de variables predeterminadas, restricciones presupuestarias agregadas)?
\item ¿Es razonable el número de condición de $\mathbf{H}_U$?
\item ¿Cambia la solución si subes $T$ a $1.5T$?
\item Recién ahora: ¿tiene sentido económico el resultado?
\end{enumerate}
\end{receta}

El orden importa. Casi todo el mundo empieza por el paso 6 --- ``esta función
impulso-respuesta se ve rara'' --- y se pasa días buscando economía donde hay un
error de signo. Los pasos 1 a 5 son mecánicos y baratos.

\section{Al escribir los resultados}

\paragraph{Muestra los jacobianos} No solo las funciones impulso-respuesta. Un
jacobiano dice mucho más sobre el mecanismo, y permite al lector comparar con
otros modelos.

\paragraph{Reporta las verificaciones} Que el método coincide con el directo,
que el horizonte es suficiente, que el estado estacionario es interior. Son tres
líneas y cambian por completo la credibilidad del trabajo numérico.

\paragraph{Distingue lo que es resultado del modelo de lo que es calibración}
Si un resultado depende de haber fijado $\tau_k$ en una grilla supuesta, eso no
es un hallazgo, es una consecuencia del supuesto. Decirlo explícitamente protege
el trabajo mucho más que esconderlo.
